% !TeX root = main.tex
%% Appendix.

%% Appendix-only table of contents.  The main body has already written its
%% entries to main.toc, so instead of reusing that file we mirror every
%% contents line produced from here on into a second file (main.atoc) and
%% typeset that one.  Everything stays local to this file.
\makeatletter
\let\apx@addcontentsline\addcontentsline
\renewcommand{\addcontentsline}[3]{%
  \apx@addcontentsline{#1}{#2}{#3}%
  \def\apx@type{#1}\def\apx@toc{toc}%
  \ifx\apx@type\apx@toc \apx@addcontentsline{atoc}{#2}{#3}\fi
}
\section*{\contentsname}
\@starttoc{atoc}
\makeatother
\newpage

\section{Related Work}
\label{sec:related}

\paragraph{Language models for formal theorem proving.}
Language models can generate formal proof steps~\citep{polu2020generative},
with subsequent systems adding tree search and reinforcement
learning~\citep{lample2022hypertree,gloeckle2024abel}, library retrieval
~\citep{yang2023leandojo}, synthetic data~\citep{xin2024deepseekprover}, subgoal
decomposition, and test-time search
~\citep{hubert2024alphaproof,ren2025deepseekproverv2,lin2025goedelprover,
wang2025kimina}. These methods primarily optimize completion of a fixed target
under a particular prover and budget, whereas we aim to explore open-ended reasoning. 

\paragraph{Autoformalization and formal libraries at scale.}
LLM autoformalization has progressed from statement-level translation
~\citep{wu2022autoformalization} and paired datasets
~\citep{azerbayev2023proofnet,ying2024lean} to large Lean developments from
mathematical textbooks~\citep{urban2026130klinesformaltopology,
wang2026m2fautomatedformalizationmathematical,gloeckle2026automatic}.
\textsc{AutoformBot} and \Atlas further demonstrate coordination and
verification across 26 books~\citep{rammal2026atlas}. Recent work has also
begun formalizing research-level results, such as the De Giorgi--Nash--Moser
theory for elliptic PDEs~\citep{armstrong2026degiorgi}. These results establish
the scalability of formal-library construction while exposing planning
problems from missing infrastructure and long-range dependencies.

\paragraph{Automated mathematical discovery and conjecturing.}
AM and Graffiti pioneered automatic concept and conjecture generation
~\citep{lenat1977am,fajtlowicz1988graffiti}, and interestingness itself became
an explicit research problem~\citep{colton2000interestingness}.
Graffiti's successors, notably TxGraffiti~\citep{davila2024txgraffiti}, have
produced conjectures that led to multiple published papers over a decade of
human--machine collaboration~\citep{davila2025txgraffiti10years}; however,
these systems operate on graph invariants and their heuristic filters are
intrinsically tied to that domain.
\citet{bengio2024aimathematician} propose an information-theoretic view
of mathematical interestingness, suggesting that valuable frameworks have small
description length while being close to many provable statements; our
definition can be seen as a computable, formal-library-grounded instantiation
of this principle.
More recent systems jointly learn conjecturing and proof in small axiomatic
domains~\citep{poesia2024minimo} or learn interestingness objectives for theory
formation~\citep{tsoukalas2025fermat}. Proposals for autonomous discovery add
novelty assessment, result selection, and closed-loop knowledge growth
~\citep{barkeshli2026structure}. Our objective supplies such a signal at formal-
library scale, grounding intrinsic interestingness in verified proof length and
relating it to a theorem's extrinsic utility. This separates the ability to
generate valid statements from the harder problem of selecting discoveries.

\paragraph{Description length and mathematical compression.}
Kolmogorov complexity measures shortest description length
~\citep{solomonoff1964,kolmogorov1965,livitanyi2008}, while proof length and
statement description length are distinct finite-corpus proxies. Resource-bounded
information~\citep{xu2020usable,blier2018description} and recent views of
mathematics as a proof hypergraph compressed by useful abstractions
~\citep{barkeshli2026structure,aksenov2026compression} motivate our
library-grounded quantities. Interestingness compares proof and description
length, while utility measures downstream proof compression; both are relative
to the available mathematical context. They are computable, finite-corpus
proxies rather than claims about uncomputable description complexity.

\section{Training the Difficulty Model \texorpdfstring{$V_\theta$}{V-theta}}
\label{app:difficulty-model}

\subsection{Data}
\label{app:extraction}

\subsubsection{Source corpus}

From our extraction of \Mathlib, we get 113,547 declarations. For each declaration we record the fully qualified name, rendered type, and file path; the file-local definitions required to interpret the target and premises; the project-local theorem declarations referenced by the accepted proof; the source proof span and its non-blank physical line count; and the dependency metadata used to change which premises are visible.

From each theorem we build a group and are able to compute our reward within such a group. Fix a root theorem $T$, let
$P$ be the premises currently visible to the model, and let $c$ be the
materialized computational-cost of that pair: the number of proof lines needed
to derive $T$ from $P$. Write $D$ for the file-local definitions shipped with
the prompt and $\Deps(\cdot)$, for the extracted dependencies of a declaration. Sample a cited lemma uniformly from the
theorem declarations among the visible premises,
\[
  L \sim \mathrm{Unif}\bigl(P)
\]
Let $\ell$ be the proof cost of $L$. Now expand it: replacing $L$ by its
own dependencies gives the expanded premise set $P^{-}$, and dropping each of
its premises independently gives the thinned premise set $P^{.8}$,
\[
  P^{-} = \bigl(P \setminus \{L\}\bigr) \cup \Deps(L),
  \qquad
  P^{.8} = \{\, p \in P^{-} : u_p < 0.8 \,\},
  \quad u_p \sim \mathrm{Unif}[0,1].
\]
A proof from $P^{-}$ can no longer cite $L$ and must derive it, so it costs
$c+\ell$. Writing $D_{\rm pre}$ for the definitions visible before the
expansion and $D(L)$ for those of $L$ alone, the edge emits six prompts
$(\text{target}\mid\text{definitions},\,\text{premises})$ together with their
observed computational-cost labels.

\begin{center}
\small
\begin{tabular}{@{}llp{0.42\linewidth}@{}}
\toprule
Role & Prompt $\mapsto$ Label & Meaning \\
\midrule
\role{base}             & $(T \mid D,\, P^{-}) \mapsto c+\ell$ & Predict the difficulty of proving $T$ when $L$ is not among the given premises; one would have to prove $L$ along the way, so the answer grows by its $\ell$ lines. \\
\role{remaining\_keep}  & $(T \mid D,\, P^{-}\cup\{L\}) \mapsto c$ & Predict the difficulty of proving $T$ when $L$ is handed back as a premise; it can be cited directly, so the answer returns to $c$. \\
\role{remaining\_remove}& $(T \mid D_{\rm pre},\, P) \mapsto c$ & Predict the difficulty of proving $T$ from the premises available before $L$ was expanded; the answer is again $c$, reached from a different set of premises. \\
\role{lemma\_context}   & $(L \mid D,\, P^{-}) \mapsto \ell$ & Predict the difficulty of proving $L$ itself from the premises the \role{base} row was given; this is exactly the $\ell$ lines that \role{base} had to add. \\
\role{lemma\_native}    & $(L \mid D(L),\, \Deps(L)) \mapsto \ell$ & Predict the difficulty of proving the same $L$ from the premises it was actually proved from in \Mathlib; the answer must again be $\ell$. \\
\role{base\_drop}       & $(T \mid D,\, P^{.8}) \mapsto \varnothing$ & Predict the difficulty of proving $T$ after each premise is dropped at random; it carries no label, since removing premises cannot make a proof shorter, and feeds only the one-sided penalty. \\
\bottomrule
\end{tabular}
\end{center}

All six rows remain adjacent during reward construction, so that the reward can be evaluated within a group without any cross-batch bookkeeping.  Roots, not individual prompt rows, determine the data
split, so variants of one theorem never straddle the train/validation boundary.

For each root, generation starts from $P=\Deps(T)$ and
$c=\Nlines(T)$ and runs three deterministically seeded random expansion walks.
After an edge is emitted, the walk continues from $(P^-,c+\ell)$. Each walk
reservoir-samples at most three eligible edges, stops after 300 expansion
attempts or once the visible premise set reaches 150 premises, and excludes prompts above
60,000 characters. The walks and reservoirs use seed 31; validation roots are
randomly selected with probability $p=0.01$.

The final dataset contains 18,318 training groups (109,908 prompts) and 923
validation groups (5,538 prompts).  Grounded labels have mean 92.3, median 49,
90th percentile 245, and maximum 500 proof lines.  

\subsubsection{Prompt}

Every model receives the same instruction template.  The frontier-model
evaluations permit free-form reasoning but require the final integer after
\texttt{\#\#\#\#}; training uses identical semantic content.

\begin{lstlisting}[basicstyle=\ttfamily\footnotesize,frame=single]
You are given Lean definitions, premises, and a target theorem statement.
Estimate the proof difficulty: the total number of proof lines required to
derive the target theorem from these premises.

The definitions are file-local definitions that the statement and premises
depend on. Use them only to understand what the statement means -- do NOT
judge difficulty based on them, and do not count them as proof lines.

Pay attention to exactly which premises are available. Do not infer
difficulty from the number of premises alone.

Do not draft the proof. Just estimate the difficulty from the dependency
structure. Give a concise estimate, then put the final answer on its own line:

#### <integer>

Definitions:
[numbered Lean declarations]

Premises:
[numbered Lean hypotheses]

Target theorem:
theorem target : [rendered type] := by

What is the total proof difficulty (number of proof lines)?
\end{lstlisting}

Hypothesis names are normalized positionally to \texttt{h1}, \texttt{h2}, and so
on; the target is renamed \texttt{target}; and any proof body after
\texttt{:= by} is stripped.  The model sees file-local definitions first, then
premises, then the target.

\subsection{Model and Optimization}
\label{app:optimization}

\subsubsection{Training objective}

Write $m(x)=\lgp{x}$, let $v_{rj}$ be sampled completion $j$ for role $r$,
and let $\bar v_r$ be the mean of the parsed completions for that role within
the batch; we abbreviate the six role means of \role{base},
\role{remaining\_keep}, \role{remaining\_remove}, \role{lemma\_context},
\role{lemma\_native}, and \role{base\_drop} as $\bar v_{\rm base}$,
$\bar v_{\rm keep}$, $\bar v_{\rm rem}$, $\bar v_{\rm ctx}$,
$\bar v_{\rm nat}$, and $\bar v_{\rm drop}$. Algorithm~\ref{alg:reward} gives
the per-completion reward used before GRPO normalization. Define the
Bellman residual $a(x;y,z)=|m(x)-m(y+z)|$ and the agreement residual
$q(x,y)=|m(x)-m(y)|$.

\begin{algorithm}[t]
\caption{\textsc{TrainingReward} --- exact per-completion reward.}
\label{alg:reward}
\begin{algorithmic}[1]
\Require completion $v=v_{rj}$ of role $r$; role means $\bar v_{\rm base}$,
  $\bar v_{\rm keep}$, $\bar v_{\rm rem}$, $\bar v_{\rm ctx}$,
  $\bar v_{\rm nat}$, $\bar v_{\rm drop}$; label $c_r$
\Ensure reward $\mathcal{R}_{rj}$
\If{$v$ does not contain exactly one \texttt{</think>} followed by exactly one
  terminal \texttt{\#\#\#\# non-negative-integer}}
  \State \Return $-2000$
\EndIf
\State $e_{\rm truth}\gets 0$ if $r=\role{base\_drop}$, else $q(v,c_r)$
\If{$r=\role{base}$}
  \State $e_{\rm add}\gets\tfrac14\bigl[
    a(v;\bar v_{\rm keep},\bar v_{\rm ctx})+a(v;\bar v_{\rm keep},\bar v_{\rm nat})$
  \Statex \hspace{5.9em}$
    +\,a(v;\bar v_{\rm rem},\bar v_{\rm ctx})+a(v;\bar v_{\rm rem},\bar v_{\rm nat})\bigr]$
  \State $e_{\rm drop}\gets[\,m(v)-m(\bar v_{\rm drop})\,]_+$
\ElsIf{$r=\role{remaining\_keep}$ or $r=\role{remaining\_remove}$}
  \State $e_{\rm add}\gets\tfrac12\bigl[
    a(\bar v_{\rm base};v,\bar v_{\rm ctx})+a(\bar v_{\rm base};v,\bar v_{\rm nat})
    \bigr];\quad e_{\rm drop}\gets0$
\ElsIf{$r=\role{lemma\_context}$}
  \State $e_{\rm add}\gets\tfrac13\bigl[
    a(\bar v_{\rm base};\bar v_{\rm keep},v)+a(\bar v_{\rm base};\bar v_{\rm rem},v)
    +q(v,\bar v_{\rm nat})\bigr];\quad e_{\rm drop}\gets0$
\ElsIf{$r=\role{lemma\_native}$}
  \State $e_{\rm add}\gets\tfrac13\bigl[
    a(\bar v_{\rm base};\bar v_{\rm keep},v)+a(\bar v_{\rm base};\bar v_{\rm rem},v)
    +q(v,\bar v_{\rm ctx})\bigr];\quad e_{\rm drop}\gets0$
\Else \Comment{$r=\role{base\_drop}$}
  \State $e_{\rm add}\gets0;\quad
    e_{\rm drop}\gets[\,m(\bar v_{\rm base})-m(v)\,]_+$
\EndIf
\State \Return $-0.50e_{\rm truth}-0.35e_{\rm add}-0.15e_{\rm drop}$
\end{algorithmic}
\end{algorithm}

\subsubsection{Hyperparameters}

\begin{center}
\small
\begin{tabular}{@{}ll@{\qquad}ll@{}}
\toprule
Parameter & Value & Parameter & Value \\
\midrule
Base model & Qwen3.6-27B & Global batch & 384 samples \\
Rollout batch & 48 prompts & Samples per prompt & 8 \\
Rollout temperature & 1.0 & Context limit & 131,072 \\
Prompt limit & 65,536 & Response limit & 8,192 \\
Optimizer & Adam & Learning rate & $10^{-6}$, constant \\
$\beta_1,\beta_2$ & $0.9,\,0.98$ & Weight decay & 0.1 \\
GRPO clip & $0.20\,/\,0.28$ & KL coefficient & 0.01 \\
Steps & 350 & Policy loss & per token \\
\bottomrule
\end{tabular}
\end{center}

\subsection{Evaluation Protocols}
\label{app:evaluation}

Evaluation is performed with greedy decoding and an 8,192-token response cap, to match training.
MAE, median absolute error, log-space $R^2$,
and Spearman $\rho$ are computed over all valid outputs. Every row from Claude and the trained model $V_\theta$ is successfully parsed, whereas GPT-5.5 produced 6 invalid outputs across all 4,615 prompts.

Calibration points in Figure~\ref{fig:difficulty-results} divide $[0,450]$ proof lines into five equal-width 90-line bins.
Within each observed-difficulty bin the marker gives the mean observed label against the mean prediction, and the shaded band spans the 25th--75th percentiles of predictions. 
Binning is used only for the plot on the left of Figure~\ref{fig:difficulty-results}; every aggregate metric on the right is computed on unbinned rows.

\section{Training a Conjecturing Model}
\label{app:conjecturing}

\subsection{Conjecturing Model Training}
\label{app:forward-training}

\subsubsection{Data construction}

The theorem-proposal data are reconstructed from the \role{base\_drop} rows of the corpus from \ref{app:extraction}, and only the premises are used here.
Premises are renamed positionally, and contexts with fewer than 16 premises or more than 12,288 tokens are removed. 
From this we randomly select 10,000 sets of premises for training and an additional 512 for validation, each set containing a median of 77 premises.
The model sees definitions and premise statements but never the source theorem or its proof. 
Section~\ref{app:proposal-prompt} gives the verbatim system and user templates.

\subsubsection{Compilation, relevance, and triviality}

A response must contain exactly one nonempty expression between
\texttt{<<<STATEMENT>>>} and \texttt{<<<END>>>}. Pantograph~\citep{aniva2025pantograph} compiles the expression as the type of a standalone theorem against the \Mathlib
environment, to ensure it is syntactically valid. We only allow for the model to produce theorems and reject declarations, definitions, proof commands,
\texttt{:=}, and \texttt{sorry}; parse failures receive $-0.5$ and expressions
that fail compilations receive $-0.25$.

For a compiled type, the global definition graph computes
$L(T\mid P)$ by subtracting the vocabulary exposed by the supplied premises.
The proposal must share at least one non-generic identifier with a
premise/definition context; failure receives $-0.25$. Direct premise restatements (as measured by token-level similarity) are marked trivial. Other candidates are tested for triviality using \texttt{assumption}, \texttt{rfl}, \texttt{simp}, \texttt{tauto},
and \texttt{simp\_all}; a candidate closed by any one receives reward of $0$. Only valid,
relevant, nontrivial statements are sent to the $V_\theta$ LLM to judge the difficulty, and they then receive reward according to Equation~\ref{eq:conjecture-reward}.

\subsubsection{Optimization configuration}

\begin{center}
\small
\begin{tabular}{@{}ll@{\qquad}ll@{}}
\toprule
Parameter & Value & Parameter & Value \\
\midrule
Initialization & Qwen3.6-27B & Steps & 75 \\
Samples per prompt & 8 & Rollout temperature & 1.0 \\
Global / rollout batch & 192 / 24 & Context / prompt limit & 12,800 / 12,288 \\
Response limit & 512 & Planned updates & 100 \\
Optimizer & Adam & Learning rate & $10^{-6}$, constant \\
$\beta_1,\beta_2$ & $0.9,0.98$ & Weight decay & 0.1 \\
GRPO clip & $0.20/0.28$ & KL penalty & $0$ \\
Save / evaluation interval & 5 / 25 & Policy loss & per token \\
\bottomrule
\end{tabular}
\end{center}

\subsubsection{Theorem-proposal prompt}

Bracketed fields such as \texttt{[NUMBERED PREMISES]} are populated mechanically
with the corresponding Lean declarations; they are not additional instructions.

% !TeX root = main.tex
%% Body of the theorem-proposal prompt (no heading: the caller supplies one,
%% at section level in the short version and subsubsection level in the full
%% appendix).
\label{app:proposal-prompt}

The system message is:

\begin{lstlisting}[basicstyle=\ttfamily\scriptsize,frame=single,breaklines=true]
You are an expert Lean 4 and Mathlib theorem proposer. Answer immediately with the tagged theorem statement. Do not think aloud, explain, explore, or write a proof. Your first output must be <<<STATEMENT>>> and the whole response should stay under 256 tokens.
\end{lstlisting}

The user message is:

\begin{lstlisting}[basicstyle=\ttfamily\scriptsize,frame=single,breaklines=true]
Given the Lean premises below, propose one new, interesting theorem statement that is a plausible logical consequence of them. The statement will be scored by estimated proof difficulty divided by its premise-relative definition length.

A strong theorem is nontrivial, concise, and uses multiple premises in a meaningful way. It must be a well-typed Lean proposition that could actually be proved from the available premises. Do not output False, an unsupported conjecture, a restatement of one premise, a vacuous implication, or a reflexive identity such as `x = x` or `P <-> P`. Use mathematical objects or operations supplied by at least two different premises. The output must be a standalone proposition: do not mention local aliases such as `h1` in the generated statement.

Prefer a short consequence obtained by instantiating and chaining the exact premises over a broad theorem with many freshly invented binders. Copy namespaces, binder types, typeclass assumptions, and function argument order exactly from the declarations below. Do not guess shorthand names or notation that is not shown. If you use a premise name such as `h2`, apply it only at its displayed signature.

Silently type-check the final expression before emitting it: fully apply maps that need arguments, give every new variable an explicit type, and use each operation only on values of its declared type (for example, do not add or subtract propositions). The output is a type, so do not use `let`, `by`, `:=`, holes, tactics, proof terms, or `#` cardinality notation. Do not print this check.

# Definitions for understanding the premise vocabulary
[NUMBERED DEFINITIONS, WHEN PRESENT]

# Available premise declarations ([POSITIONAL PREMISE NAMES])
[NUMBERED PREMISES]

# Output contract
Output only a single standalone Lean type expression between the tags. Do not reference the premise aliases above. If new variables are necessary, bind them with a leading `forall`/`forall`. Do not include a theorem name, `:= by`, proof, code fence, or commentary.

<<<STATEMENT>>>
<one concise Lean theorem statement>
<<<END>>>
\end{lstlisting}



\subsection{Conjecturing Model Evaluation}
\label{app:forward-eval}

\subsubsection{Balanced verified cohort}

The eight strata we select premises from are Algebra, Analysis, Number Theory, Measure/Probability, Geometry/Topology, Combinatorics, Category/Algebraic Geometry, and Foundations.
For each group, we sample theorems from each of the three models we evaluate (Claude, the trained conjecturing model, and the base model), and pass it to a Claude Code with Claude Opus 4.6 as a proof agent. We repeat until we have 20 distinct theorems from each area.

The proving agent first attempts the statement byte-for-byte. If that fails, a separate
marginal-repair pass may add a necessary hypothesis, correct an index or bound,
or narrowly weaken a conclusion. A repair is accepted only if it compiles,
preserves the counts of $\forall,\exists,\leftrightarrow,\land,\lor,\neg$, passes
anti-vacuity checks, and has a token-set Jaccard similarity at least $0.70$, to disallow the statement becoming irrelevant or too easy.
We recompute proof lines and conditional length from the final compiled declaration and compiled proof, so the reported score is ground-truth interestingness under the proving pipeline rather
than the difficulty model's training reward.

\subsubsection{Mathlib containment}
\label{app:containment}

Claude Opus 4.6 at temperature 0
independently judges all 480 verified statements with proposer, area, and model
hidden. Each prompt includes the final statement, verified proof body, and
referenced \Mathlib declarations as non-exhaustive evidence. The proof scaffold
is pinned to \Mathlib revision
\texttt{905b95...}. Scores range from 1 (not contained) to 5 (an exact or trivially equivalent named
declaration exists); score 4 denotes a few-line wrapper or immediate
specialization. The judge must name a concrete declaration for score 5 and
returns structured JSON. Figure~\ref{fig:mathlib-containment} reports all of the recorded
statements per model.  The complete rubric and judge prompt are provided in
Section~\ref{app:containment-prompt}.

\subsubsection{Mathlib-containment rubric and judge prompt}

% !TeX root = main.tex
%% Body of the Mathlib-containment rubric and judge prompt (no heading: the
%% caller supplies one).
\label{app:containment-prompt}

The five-level rubric is:

\begin{lstlisting}[basicstyle=\ttfamily\scriptsize,frame=single,breaklines=true]
5 -- Fully contained. The exact statement (or a direct, trivially equivalent reformulation) exists in Mathlib as a named declaration. A user could cite it directly, or with a thin wrapper.

4 -- Substantially contained. The statement is provable in a few lines from existing Mathlib lemmas, or a strictly more general version is in Mathlib and the target is an immediate specialization.

3 -- Partially contained. Core ingredients (definitions and key supporting lemmas) exist in Mathlib, but the headline statement itself is not there and would require nontrivial assembly.

2 -- Minimally contained. Only background definitions or unrelated prerequisites are available; the substantive content is missing.

1 -- Not contained. Mathlib does not contain the statement or the necessary specialized definitions.

Operational guidance:
- Judge library containment, not mathematical validity, beauty, novelty, importance, or apparent difficulty.
- "A few lines" means a short Lean proof assembled directly from existing declarations, not merely a short informal mathematical argument.
- Use score 5 only when there is a concrete matching named declaration (possibly with symmetry, simplification, notation changes, or an immediate specialization). Identify that declaration whenever possible.
- A verified proof and its referenced Mathlib declarations are supplied as evidence. They are not exhaustive: a generated proof may overlook a more direct existing theorem, and proof length is not by itself dispositive.
- Do not infer a high score merely because the theorem is elementary. Conversely, do not infer a low score merely because its generated proof is long.
- If the exact declaration name is uncertain, report that uncertainty and avoid score 5 unless the API match is otherwise unmistakable.
\end{lstlisting}

Each item is judged with:

\begin{lstlisting}[basicstyle=\ttfamily\scriptsize,frame=single,breaklines=true]
You are an impartial expert judge of Lean 4 Mathlib library containment.

Apply the supplied rubric exactly. Estimate how much of the target theorem is already contained in Mathlib. The theorem has been verified, so do not judge whether it is true. Do not judge its mathematical novelty, importance, elegance, or difficulty.

The proof body and referenced declarations are supporting evidence, not an exhaustive Mathlib search. A generated proof can miss a direct existing theorem. Treat score 5 as requiring a concrete matching named declaration or an unmistakably direct API match. Distinguish a genuinely few-line Mathlib wrapper (score 4) from nontrivial assembly of ingredients (score 3).

<rubric>
[RUBRIC ABOVE]
</rubric>

<target_theorem>
[STATEMENT]
</target_theorem>

<verified_proof_body lines="[LINE COUNT]">
[PROOF BODY]
</verified_proof_body>

<referenced_mathlib_declarations>
[NAMES AND TYPES, OR NO RECOVERED LIST]
</referenced_mathlib_declarations>

Return exactly one JSON object and no Markdown or surrounding commentary:
{
  "score": <integer 1-5>,
  "matching_declaration": "<concrete Mathlib declaration for score 5, best candidate otherwise, or none>",
  "supporting_declarations": ["<zero or more important declaration names>"],
  "rationale": "<concise explanation tied to the rubric>",
  "confidence": "<high | medium | low>"
}
\end{lstlisting}


\newcommand{\genexheading}{\subsection}
%% Heading level is chosen by the caller (\section in the short version,
%% \subsection in the full appendix).
\providecommand{\genexheading}{\section}
\genexheading{Selected Generations}
\label{app:generated-examples}

Each item gives the final compiled Lean statement followed by a mathematical
informalization. 

\begin{list}{}{%
  \setlength{\leftmargin}{0pt}%
  \setlength{\rightmargin}{0pt}%
  \setlength{\labelwidth}{0pt}%
  \setlength{\labelsep}{0pt}%
  \setlength{\itemindent}{0pt}%
  \setlength{\itemsep}{1.25em}%
  \setlength{\parsep}{0pt}%
  \setlength{\topsep}{0.75em}%
}

\item
\begin{samepage}
\begin{lstlisting}[basicstyle=\ttfamily\footnotesize,frame=single,breaklines=true]
∀ q : ℍ,
  q * q = -1 ↔ q.re = 0 ∧ ‖q‖ = 1
\end{lstlisting}
\noindent\textbf{Quaternionic square roots of minus one.}
A real quaternion squares to $-1$ if and only if its real component is zero
and its norm is one. Equivalently, the square roots of $-1$ are precisely the
unit purely imaginary quaternions.
\end{samepage}

\begin{comment}\item
\begin{samepage}
\begin{lstlisting}[basicstyle=\ttfamily\footnotesize,frame=single,breaklines=true]
∀ (α : Type _) [Group α], ∀ a b : α,
  a * b = 1 ↔
    b * a = 1 ∧
      ∀ c : α, c * a = 1 ↔ a * c = 1
\end{lstlisting}
\noindent\textbf{One-sided inverses in a group.}
For elements $a,b$ of a group, $ab=1$ exactly when $ba=1$ and, for every
$c$, the equation $ca=1$ is equivalent to $ac=1$. All of these equations say
that the other element is the unique inverse of $a$.
\end{samepage}
\end{comment}

\begin{comment}
\item
\begin{samepage}
\begin{lstlisting}[basicstyle=\ttfamily\footnotesize,frame=single,breaklines=true]
∀ x : ℝ,
  0 < x ↔
    Real.log (Real.exp x + 1) > Real.log 2
\end{lstlisting}
\noindent\textbf{Positivity through a logarithmic transform.}
A real number $x$ is positive if and only if
$\log(1+e^x)>\log 2$. Thus the softplus transform crosses its value at zero
exactly when $x$ crosses zero.
\end{samepage}
\end{comment}

\item
\begin{samepage}
\begin{lstlisting}[basicstyle=\ttfamily\footnotesize,frame=single,breaklines=true]
∀ z : UpperHalfPlane,
  ∃ n : ℕ, 0 < n ∧
    ‖Complex.exp
      (2 * ↑Real.pi * Complex.I * ↑z) ^ n‖ < 1 / 2
\end{lstlisting}
\noindent\textbf{Decay of the upper-half-plane exponential.}
For every $z$ in the complex upper half-plane, some positive power of
$e^{2\pi i z}$ has absolute value less than $1/2$. 
\end{samepage}

\begin{comment}
\item
\begin{samepage}
\begin{lstlisting}[basicstyle=\ttfamily\footnotesize,frame=single,breaklines=true]
∀ r : ℝ, r < 1 →
  (0 < r ∧ r < 1 ↔
    Real.log ((1 + r) / (1 - r)) > 0)
\end{lstlisting}
\noindent\textbf{Positivity of a logarithmic ratio.}
For a real $r<1$, the quantity
$\log\!\left((1+r)/(1-r)\right)$ is positive exactly when $0<r<1$.
\end{samepage}
\end{comment}

\item
\begin{samepage}
\begin{lstlisting}[basicstyle=\ttfamily\footnotesize,frame=single,breaklines=true]
∀ (n : Type _) [Fintype n] [DecidableEq n]
    (R : Type _) [CommRing R]
    (A : Matrix n n R),
  A * A = 1 ↔
    A.det * A.det = 1 ∧
      A.adjugate = A.det • A
\end{lstlisting}
\noindent\textbf{Involutory matrices and their adjugates.}
For a finite square matrix $A$ over a commutative ring, $A^2=I$ if and only if
$(\det A)^2=1$ and $\operatorname{adj}(A)=(\det A)A$.
\end{samepage}

\item
\begin{samepage}
\begin{lstlisting}[basicstyle=\ttfamily\footnotesize,frame=single,breaklines=true]
∀ κ : Cardinal,
  κ * κ = κ ↔
    κ = 0 ∨ κ = 1 ∨ Cardinal.aleph0 ≤ κ
\end{lstlisting}
\noindent\textbf{Multiplicatively idempotent cardinals.}
A cardinal is unchanged by squaring exactly when it is $0$, $1$, or infinite
(equivalently, at least $\aleph_0$).
\end{samepage}

\item
\begin{samepage}
\begin{lstlisting}[basicstyle=\ttfamily\footnotesize,frame=single,breaklines=true]
∀ z : UpperHalfPlane,
  ∃ w : UpperHalfPlane,
    w = -(1 : ℂ) / (z : ℂ) ∧
      ∀ n : ℕ, n > 0 ↔
        (↑n : ℂ) * (z : ℂ) ≠ 0
\end{lstlisting}
\noindent\textbf{Modular inversion in the upper half-plane.}
For every upper-half-plane point $z$, the modular inverse $w=-1/z$ is again in
the upper half-plane. Moreover, for every natural number $n$, positivity of
$n$ is equivalent to the nonvanishing of the complex product $nz$.
\end{samepage}

\item
\begin{samepage}
\begin{lstlisting}[basicstyle=\ttfamily\footnotesize,frame=single,breaklines=true]
∀ (R : Type _) [CommRing R],
  ∀ (W' : WeierstrassCurve.Jacobian R)
      (P : Fin 3 → R),
    W'.Equation (W'.neg P) ↔ W'.Equation P
\end{lstlisting}
\noindent\textbf{Negation preserves the Jacobian equation.}
For a Jacobian Weierstrass model over any commutative ring, negating a
projective point preserves the curve equation: $-P$ lies on the curve exactly
when $P$ does.
\end{samepage}

\item
\begin{samepage}
\begin{lstlisting}[basicstyle=\ttfamily\footnotesize,frame=single,breaklines=true]
∀ (K : Type _) [Field K] (f : Polynomial K),
  Function.Injective
    (fun g : Polynomial K => g * f) ↔
      f ≠ 0
\end{lstlisting}
\noindent\textbf{Injectivity of polynomial multiplication.}
Over a field, right multiplication by a polynomial $f$ is injective on the
polynomial ring if and only if $f$ is nonzero.
\end{samepage}

\item
\begin{samepage}
\begin{lstlisting}[basicstyle=\ttfamily\footnotesize,frame=single,breaklines=true]
∀ (K : Type _) [Field K] [NumberField K]
    (x : K),
  x ≠ 0 ↔
    ∀ w : NumberField.InfinitePlace K,
      (NumberField.InfinitePlace.mk
        (NumberField.InfinitePlace.embedding w)) x > 0
\end{lstlisting}
\noindent\textbf{Positivity at every infinite place.}
An element of a number field is nonzero exactly when its value at every
infinite place is strictly positive, where an infinite place evaluates through
the absolute value associated with its embedding.
\end{samepage}

\item
\begin{samepage}
\begin{lstlisting}[basicstyle=\ttfamily\footnotesize,frame=single,breaklines=true]
∀ (α : Type _) [GeneralizedBooleanAlgebra α]
    (a b : α),
  a = b ↔
    ∀ c : α, a \ c = b \ c ∧ c \ a = c \ b
\end{lstlisting}
\noindent\textbf{Equality via relative complements.}
Two elements $a,b$ of a generalized Boolean algebra are equal exactly when,
for every $c$, removing $c$ from them gives the same result and removing them
from $c$ also gives the same result.
\end{samepage}

\item
\begin{samepage}
\begin{lstlisting}[basicstyle=\ttfamily\footnotesize,frame=single,breaklines=true]
∀ (V : Type _) [Finite V] (G : SimpleGraph V),
  G.edgeSet = ∅ ↔
    ∀ s : Set V, G.IsClique s ↔ s.ncard < 2
\end{lstlisting}
\noindent\textbf{Empty graphs characterized by their cliques.}
A finite simple graph has no edges if and only if its cliques are precisely the
vertex sets containing fewer than two vertices.
\end{samepage}

\item
\begin{samepage}
\begin{lstlisting}[basicstyle=\ttfamily\footnotesize,frame=single,breaklines=true]
∀ (V : Type _) (G : SimpleGraph V) [Finite V],
  G.IsClique (Set.univ : Set V) ↔
    ∀ (M : Gᶜ.Subgraph),
      M.IsMatching ↔ M.verts = ∅
\end{lstlisting}
\noindent\textbf{Completeness through matchings in the complement.}
A finite simple graph is complete if and only if a subgraph of its complement
is a matching exactly when that subgraph has no vertices.
\end{samepage}

\item
\begin{samepage}
\begin{lstlisting}[basicstyle=\ttfamily\footnotesize,frame=single,breaklines=true]
∀ n : ℕ, n > 0 ↔
  ∀ f : Polynomial ℝ, f.natDegree ≤ n - 1 →
    ((∀ k < n,
      Polynomial.eval
        (Polynomial.Chebyshev.node n k) f = 0) ↔
      f = 0)
\end{lstlisting}
\noindent\textbf{Vanishing at Chebyshev nodes.}
A natural number $n$ is positive exactly when every real polynomial of degree
at most $n-1$ that vanishes at all $n$ Chebyshev nodes is the zero polynomial
(and the zero polynomial vanishes at those nodes).
\end{samepage}

\item
\begin{samepage}
\begin{lstlisting}[basicstyle=\ttfamily\footnotesize,frame=single,breaklines=true]
∀ (C : Type _) [CategoryTheory.Category C]
    [CategoryTheory.Abelian C]
    (X Y : C) (f : X ⟶ Y),
  CategoryTheory.Mono f →
    (CategoryTheory.IsIso f ↔
      ∀ (Z : C) (g : Y ⟶ Z),
        (g = 0 ↔ f ≫ g = 0))
\end{lstlisting}
\noindent\textbf{Detecting isomorphisms in an abelian category.}
Let $f:X\to Y$ be a monomorphism in an abelian category. Then $f$ is an
isomorphism exactly when precomposition with $f$ detects zero morphisms out of
$Y$: for every $g:Y\to Z$, one has $g=0$ if and only if $g\circ f=0$.
\end{samepage}

\end{list}



\section{Inference-Time Pruning Ablation}
\label{app:inference-pruning}

In this experiment, we aim to isolate the effect of the promotion rule in the recursive discovery loop of
Section~\ref{sec:fwd}. All four trajectories we study start from the set of 80 premises $P_0$ from graph theory, and share the same Claude Opus 4.6 conjecturer model, and Claude Code based proving procedure. 
Each initial generation pass proposes 400 statements. Every prompt contains 20 premises: five sampled from the immediately preceding set $P_n \backslash P_{n-1}$ and 15 from the initial pool and older promoted rounds $P_{n-1}$.

These statements are proved independently in Lean. We compare four post-verification promotion rules: \emph{no pruning} retains every verified novel statement; \emph{random} retains ten uniformly at random; \emph{proof length} retains the ten statements with the longest verified proofs; and \emph{interestingness} retains the ten statements with the largest ratio
\begin{equation}
  \widehat I_{\rm ver}(T\mid P)
    =100\,\frac{\text{lines in the accepted proof of }T}
                    {\text{characters in the normalized statement of }T}.
  \label{eq:verified-forward-interestingness}
\end{equation}

\subsection{Comparison across promotion policies}

All policies share the same initial premises and differ only in which verified statements are promoted. Table~\ref{tab:forward-pruning-primary} compares the statements promoted over
the ten rounds. Mean interestingness is computed within each round and then
averaged across rounds, giving every round equal weight.

\begin{table}[H]
  \centering
  \small
  \begin{tabular}{@{}lrrrr@{}}
    \toprule
    Promotion rule  & Mean $\widehat I_{\rm ver}$ & Median $\widehat I_{\rm ver}$ & Mean proof lines \\
    \midrule
    No pruning       & 1.313 & 1.114 & 3.80 \\
    Random           & 1.556 & 1.175 & 4.39 \\
    Proof length     & 3.259 & 2.480 & \textbf{11.48} \\
    Interestingness  & \textbf{3.899} & \textbf{3.380} & 7.41 \\
    \bottomrule
  \end{tabular}
  \caption{\textbf{Quality of the promoted statements}}
  \label{tab:forward-pruning-primary}
\end{table}

Unsurprisingly, interestingness pruning produces the highest mean and median interestingness,
while proof-length pruning produces the longest proofs.
A direct advantage on promoted interestingness is partly expected because one
policy selects on that quantity. We therefore separately evaluate every unique
candidate successfully proved in Lean \emph{before} the current round's
promotion rule is applied. Table~\ref{tab:forward-pruning-candidates} reports
their number, mean and median interestingness, and mean proof length across all
ten rounds, using the same aggregation as
Table~\ref{tab:forward-pruning-primary}.

\begin{table}[H]
  \centering
  \small
  \begin{tabular}{@{}lrrrr@{}}
    \toprule
    Promotion rule  & Mean $\widehat I_{\rm ver}$ & Median $\widehat I_{\rm ver}$ & Mean proof lines \\
    \midrule
    No pruning      & 1.313 & 1.114 & 3.80 \\
    Random          & 1.461 & 1.271 & 3.74 \\
    Proof length    & 1.694 & 1.167 & \textbf{5.46} \\
    Interestingness & \textbf{2.332} & \textbf{1.754} & 4.56 \\
    \bottomrule
  \end{tabular}
  \caption{\textbf{Verified candidates before selection.}}
  \label{tab:forward-pruning-candidates}
\end{table}


The judge sees four statements at a time, one from each policy, and ranks them by mathematical interestingness. Policy names, proofs, and scores are hidden. We evaluate 100 such groups, and show the results below.

\begin{table}[H]
  \centering
  \small
  \begin{tabular}{@{}lrrr@{}}
    \toprule
    Promotion rule & Ranked first & Ranked in top two & Mean rank $\downarrow$ \\
    \midrule
    No pruning      & 13\% & 34\% & 2.89 \\
    Random          & 8\%  & 29\% & 2.95 \\
    Proof length    & 19\% & 64\% & 2.36 \\
    Interestingness & \textbf{60\%} & \textbf{73\%} & \textbf{1.80} \\
    \bottomrule
  \end{tabular}
  \caption{\textbf{Direct policy-blinded judgment of mathematical
  interestingness.}}
  \label{tab:forward-pruning-interest-judge}
\end{table}

The judge also evaluates ten groups of ten statements promoted by each policy in each round as a group. Policy names and scores are hidden. The judge rates the group's overall quality and diversity and identifies repeated theorem families.
For no pruning, which retains more than ten statements, we select ten at random.

\begin{table}[H]
  \centering
  \small
  \begin{tabular}{@{}lrrrr@{}}
    \toprule
    Promotion rule & Quality (1--5) & Diversity (1--5) & Families per ten & Redundant \\
    \midrule
    No pruning      & $2.50\pm0.53$ & $1.30\pm0.48$ & $2.60\pm0.97$ & $50.0\%\pm26.7\%$ \\
    Random          & $2.40\pm0.52$ & $1.20\pm0.42$ & $2.90\pm0.99$ & $53.0\%\pm18.9\%$ \\
    Proof length    & $2.60\pm0.52$ & $1.30\pm0.48$ & $2.90\pm0.74$ & $50.0\%\pm18.3\%$ \\
    Interestingness & $\mathbf{3.40\pm0.52}$ & $\mathbf{2.10\pm0.32}$ & $\mathbf{3.50\pm0.71}$ & $\mathbf{28.0\%\pm7.9\%}$ \\
    \bottomrule
  \end{tabular}
  \caption{\textbf{Policy-blinded cohort judgment.} Entries are mean $\pm$
  sample standard deviation across the ten round-level cohorts for each
  promotion rule.}
  \label{tab:forward-pruning-judge}
\end{table}

\subsection{Judge prompts}

The complete judge instruction is deliberately short and uses a broad notion
of what a mathematician might appreciate:

\begin{lstlisting}[basicstyle=\ttfamily\footnotesize,breaklines=true]
You are a mathematician comparing four formally verified theorem statements.

Rank them from most to least mathematically interesting. Use a broad notion of interestingness: a mathematician might appreciate a statement because it is nontrivial, surprising, elegant, conceptually illuminating, potentially useful or reusable, connects ideas, or opens further questions. Judge the mathematical content, not notation, verbosity, statement length, or guessed proof length.

[A]
<THEOREM>

[B]
<THEOREM>

[C]
<THEOREM>

[D]
<THEOREM>

Return exactly one JSON object and no Markdown:
{"ranking":["A","B","C","D"],"rationale":"one or two sentences","confidence":"high|medium|low"}
\end{lstlisting}

The complete cohort-judge prompt is:

\begin{lstlisting}[basicstyle=\ttfamily\footnotesize,breaklines=true]
You are an impartial expert judge of a collection of ten verified Lean 4 theorem statements.

The statements are all correct. Judge the collection itself, without guessing its source or generation method.

Rate two dimensions from 1 to 5.

MATHEMATICAL QUALITY
5: Mostly nontrivial, coherent, potentially reusable results that express substantive mathematical content.
4: Generally meaningful results with some routine or overly specialized items.
3: Mixed quality; several useful statements but substantial routine conjunctions, wrappers, or narrow variants.
2: Mostly shallow, over-specialized, or mechanically assembled statements.
1: Almost entirely trivial, incoherent, vacuous, or unhelpful statements.

SEMANTIC DIVERSITY
5: Many genuinely distinct theorem families or mathematical ideas, with little redundancy.
4: Several distinct families and only limited repetition.
3: A meaningful mixture, but multiple statements repeat the same core pattern.
2: Dominated by one or two narrow families with superficial variations.
1: Nearly all statements are equivalent, nested variants, or routine repackagings.

Group statements by their substantive conclusion and proof idea, not merely surface syntax. A stronger/weaker variant or conjunction of the same core facts should normally stay in the same family.

<statements>
[1]
<THEOREM>
...
[10]
<THEOREM>
</statements>

Return exactly one JSON object and no Markdown:
{
  "quality_score": <integer 1-5>,
  "diversity_score": <integer 1-5>,
  "distinct_families": <integer 1-10>,
  "redundant_statement_indices": [<zero or more integers 1-10>],
  "family_assignments": [{"indices": [<integers>], "family": "<short label>"}],
  "rationale": "<concise explanation>",
  "confidence": "<high | medium | low>"
}
\end{lstlisting}
\clearpage

\section{Additional Figures}

\subsection{Cross-area interestingness matrix}
\label{app:interestingness}

The displayed matrix uses ten targets from each of 12 areas.  Starting from one
area-conditioned premise frontier per column, we apply three synchronous
expansion rounds: every current extracted theorem premise is replaced by all
original theorem premises used in its proof, unavailable declarations remain
terminals, and each resulting frontier is deduplicated.  Fresh $V_\theta$
predictions are obtained for all 1,440 target--condition pairs.
Each target's score is divided by its same-area score before cell medians are
formed.  Only 28 off-diagonal 95\% confidence intervals exclude one.
Across the 132 off-diagonal cell medians, the median is $0.520$ and 73.5\% are
below one. 
We see here that, as we might expect, a definition-heavy field such as measure theory becomes much less interesting when it is not given premises from analytic areas. 
We find also that some areas, such as algebra, become more interesting when supplied with premises from other fields.

\begin{figure}[H]
  \centering
  \includegraphics[width=0.8\linewidth]{figures/area_matrix.pdf}
  \caption{\textbf{Full cross-area interestingness matrix.} Each cell is the
  median target-wise ratio for the row area under premises from the column area
  versus the same-area premise condition. All premise frontiers are
  synchronously expanded through three rounds before new predictions are obtained. We select ten targets per area (120
  total); four row areas contribute nine ratios after each excludes one
  non-positive same-area score. Asterisks mark 95\% confidence intervals that exclude one.}
  \label{fig:area-matrix-full}
\end{figure}

\clearpage

\subsection{Iterative discovery graphs}

Figure \ref{fig:forward-dependency-graphs} provides more iterative discovery graphs.

\begin{figure}[H]
  \centering
  \begin{subfigure}[b]{\linewidth}
    \centering
    \includegraphics[width=\linewidth]{figures/forward_dependency_graph_algebra.pdf}
    \caption{Algebra. For monic $q$ and $r$, reducing $p r$ modulo $q r$
    equals reducing $p$ modulo $q$ and then multiplying by $r$.}
  \end{subfigure}

  \medskip
  \begin{subfigure}[b]{\linewidth}
    \centering
    \includegraphics[width=\linewidth]{figures/forward_dependency_graph_measure_probability.pdf}
    \caption{Measure/probability. If $\mu$ is absolutely continuous with
    respect to $\nu$, then under the stated measurability conditions, the sum
    of the pushforwards of $c\mu$ and $\mu$ is absolutely continuous with
    respect to $(c+1)$ times the pushforward of $\nu$.}
  \end{subfigure}

  \medskip
  \begin{subfigure}[b]{\linewidth}
    \centering
    \includegraphics[width=\linewidth]{figures/forward_dependency_graph_number_theory.pdf}
    \caption{Number theory. Applying a real invertible $2\times2$ matrix to a
    nonreal complex number and then applying its inverse restores the imaginary
    part.}
  \end{subfigure}
  \caption{\textbf{Iterative forward-discovery graphs through $P_6$.}
  Each panel shows 30 of the 240 initial premises in $P_0$, always retaining
  every initial premise used by a displayed proof. Gray lines record explicit
  dependencies in accepted Lean proofs; coral marks the ancestry of the
  theorem with the highest total $P_0$-relative interestingness among the
  statements introduced in $P_6$. Green encodes $\log_{10} I_{P_0}$ with a
  separate scale in each panel. Each subcaption reports the selected statement.}
  \label{fig:forward-dependency-graphs}
\end{figure}


